\chapter{Líneas de investigación paralelas}

Durante el desarrollo de esta investigación se realizaron dos proyectos paralelos que contribuyeron al entendimiento de los sistemas complejos y, en particular, de los autómatas celulares neuronales. Ambos estudian el espacio de reglas de los autómatas celulares, es decir, el conjunto de todas las reglas locales posibles para un número dado de estados y de vecinos, desde dos perspectivas complementarias: cómo explorarlo para encontrar reglas con un comportamiento de interés y cómo medir qué tan parecidas son dos reglas por su comportamiento. Estas preguntas se relacionan directamente con los autómatas celulares neuronales, cuyas reglas no se diseñan a mano, sino que se obtienen mediante un proceso de optimización y, por lo tanto, también deben buscarse y compararse.

El primer proyecto, \textit{Adaptive Evolution in Totalistic Cellular Automaton}, se desarrolló en la \textit{Wolfram Summer School} 2025 y se presentó en la \textit{Wolfram Technology Conference} 2025. En él se buscan reglas de autómatas celulares totalísticos que, a partir de una sola celda activa, generen formas de vida finitas, combinando filtros que reducen el espacio de reglas con algoritmos de búsqueda de inspiración biológica. El segundo, \textit{Spectral Moments with Attractor Mass: A Similarity Index for One-Dimensional Cellular Automata}, se realizó en colaboración con Sukanta Das en la \textit{Indian Summer School on Cellular Automata} 2026 y propone un índice de similitud entre reglas basado en su grafo de transiciones, aplicable incluso a reglas con distinto número de estados o de vecinos. Ambos proyectos están relacionados: el primero agrupa reglas por su comportamiento a través de un grafo de mutaciones, mientras que el segundo busca cuantificar qué tan similares son dos reglas.
